% Generated draft: requires mathematical and source review.
\documentclass[11pt]{article}
\usepackage[margin=25mm]{geometry}
\usepackage{fontspec}
\setmainfont{Latin Modern Roman}
\usepackage{amsmath,amssymb,mathtools,unicode-math}
\setmathfont{Latin Modern Math}
\usepackage{enumitem,needspace,xurl,hyperref}
\setlength{\parindent}{0pt}
\setlength{\parskip}{5pt}
\setlength{\emergencystretch}{3em}
\newcommand{\N}{\mathbb{N}}
\newcommand{\Z}{\mathbb{Z}}
\newcommand{\Q}{\mathbb{Q}}
\newcommand{\R}{\mathbb{R}}
\newcommand{\C}{\mathbb{C}}
\newcommand{\F}{\mathbb{F}}
\newcommand{\A}{\mathbb{A}}
\newcommand{\PP}{\mathbb{P}}
\newcommand{\E}{\mathbb{E}}
\newcommand{\eps}{\varepsilon}
\newcommand{\dd}{\,\mathrm{d}}
\begin{document}
\section*{UnsolvedMath ID: 30005011}
\textbf{Source problem code:} catalogue entry
\subsection{Definitions and mathematical statement}
Let \(V_1,\ldots,V_m\subseteq\C^d\) be given algebraic varieties, and let \(p_1,\ldots,p_m\in R:=\C[x_1,\ldots,x_d]\) be given polynomials. For a polynomial \(g\in R\), its restriction \(g|_{V_j}\) is the function \(V_j\to\C\) given by \(z\mapsto g(z)\).

Determine necessary and sufficient conditions on the data \((V_j,p_j)_{j=1}^m\) for the following existence statement to hold:
\[
\exists f\in\C[x_1,\ldots,x_d]\quad
\forall j\in\{1,\ldots,m\}\quad
\forall z\in V_j,\qquad f(z)=p_j(z).
\]
The requested agreement is pointwise agreement on each variety; the supplied statement imposes no derivative conditions, degree bound, disjointness assumption, or prior compatibility condition on the given polynomials.

The source does not specify whether its use of ``algebraic variety'' requires irreducibility. This convention remains unresolved, and no irreducibility hypothesis is added here. The supplied literature assessment refers to radicality hypotheses without identifying the ideals or hypotheses involved; that assessment is not adopted as an additional hypothesis or a mathematical conclusion.
\subsection{Short English statement}
Characterize when prescribed polynomial functions on finitely many algebraic varieties in \(\C^d\) are restrictions of one polynomial on the ambient space. The source's convention concerning irreducibility remains unspecified.
\subsection{Sources}
\begin{itemize}
\item Local UnsolvedMath registry: supplied record with permanent ID \(30005011\), title ``Compatible Polynomial Interpolation on Algebraic Varieties,'' and exact source problem code \texttt{catalogue entry}. No individual proposer is identified in the record; its proposed year is \(2022\).
\item Supplied report citation: \emph{Mini-Workshop: Interpolation, Approximation, and Algebra (2023)}, Oberwolfach Reports, pp.~383--428, DOI \texttt{10.4171/owr/2022/7}. \url{https://doi.org/10.4171/owr/2022/7}. The supplied background describes Oberwolfach Reports as records of mathematical workshops and their open-problem discussions. The citation's \(2023\) year and the record's \(2022\) year are retained without reconciliation.
\item Supplied literature reference: ``McKinley--Shekhtman--Tuesink report results,'' \url{https://doi.org/10.4171/owr/2022/7}. Literature assessment dated exactly \texttt{2026-08-22}: the supplied curation labels the problem partially solved and states that Hilbert--Nullstellensatz and ideal-theoretic criteria solve disjoint varieties and provide exact criteria under radicality hypotheses, while the fully general data-dependent interpolation problem remains only partially characterized. This is supplied dated curation derived from the OWR material, not an independently verified literature assessment; its time-sensitive status and cited claims require independent verification.
\item The registry labels the statement \texttt{corrected\_verified} and reports that it was checked against the official Oberwolfach report. This is a supplied verification claim, not verification performed for this draft.
\end{itemize}
\end{document}
